\chapter{Làm Lớp Vui Quá Đà}
\chapterintrobi{Có duyên là tốt, nhưng vui quá tay thì tiết học dễ trượt nhịp}{Be too funny in class}{Một vài bạn sinh ra để làm lớp cười, rồi vô tình làm bài học hụt hơi theo.}

\begin{lucbatbox}{Lục bát: Vui đến lố nhịp bài}
Em vừa chen một câu thôi\
Cả lớp cười ngặt, tiết trôi mất đà\
Nhiều khi ý tốt thật mà\
Muốn cho không khí đỡ là nặng ghê\
Chỉ là quá mức một bề\
Tiếng cười đi trước, bài về sau lưng
\end{lucbatbox}

\begin{tutuyetbox}{Thất ngôn tứ tuyệt: Duyên hơn kế hoạch}
Em rất có duyên giữa lớp này\
Nói câu nào trúng tiếng cười ngay\
Biết giữ vừa đúng lúc nữa thôi
Tiết học vừa vui chẳng chao lay
\end{tutuyetbox}

\begin{ghichu}
Làm lớp vui là một năng lực tốt. Cái cần học thêm là biết dừng đúng lúc để tiếng cười không nuốt mất trọng tâm.
\end{ghichu}

\begin{englishnote}
\textbf{too much}: quá đà.\par
\textbf{keep the balance}: giữ chừng mực.\par
\textbf{I just wanted to make everyone laugh.}: Em chỉ muốn cả lớp cười chút thôi.
\end{englishnote}